% ============================================================================
\begin{thebibliography}{99}
\small

\bibitem{cressman2009} Cressman JR, Ullah G, Ziburkus J, Schiff SJ, Barreto E. The influence of sodium and potassium dynamics on excitability, seizures, and the stability of persistent states: I. Single neuron dynamics. \emph{J Comput Neurosci} 26:159--170, 2009; Erratum, 30:781, 2011. ModelDB 136773.
\bibitem{ullah2009} Ullah G, Cressman JR, Barreto E, Schiff SJ. The influence of sodium and potassium dynamics on excitability, seizures, and the stability of persistent states: II. Network and glial dynamics. \emph{J Comput Neurosci}, 2009.
\bibitem{krishnan2011} Krishnan GP, Bazhenov M. Ionic dynamics mediate spontaneous termination of seizures and postictal depression state. \emph{J Neurosci} 31:8870, 2011.
\bibitem{jirsa2014} Jirsa VK et al. On the nature of seizure dynamics. \emph{Brain} 137:2210, 2014.
\bibitem{jirsa2017} Jirsa VK et al. The Virtual Epileptic Patient: individualized whole-brain models of epilepsy spread. \emph{NeuroImage} 145:377, 2017.
\bibitem{hashemi2020} Hashemi M et al. The Bayesian Virtual Epileptic Patient. \emph{NeuroImage} 217:116839, 2020.
\bibitem{sanzleon2013} Sanz-Leon P et al. The Virtual Brain: a simulator of primate brain network dynamics. \emph{Front Neuroinform} 7:10, 2013.
\bibitem{chizhov2018} Chizhov AV et al. Minimal model of interictal and ictal discharges ``Epileptor-2''. \emph{PLoS Comput Biol}, 2018.
\bibitem{izhikevich2003} Izhikevich EM. Simple model of spiking neurons. \emph{IEEE Trans Neural Netw} 14:1569, 2003.
\bibitem{compte2003} Compte A et al. Cellular and network mechanisms of slow oscillatory activity. \emph{J Neurophysiol} 89:2707, 2003.
\bibitem{doyon2016} Doyon N et al. Chloride regulation: a dynamic equilibrium crucial for synaptic inhibition. \emph{Neuron} 89:1157, 2016.
\bibitem{gulledge2013} Gulledge AT et al. A sodium-pump-mediated afterhyperpolarization in pyramidal neurons. \emph{J Neurosci} 33:13025, 2013.
\bibitem{during1992} During MJ, Spencer DD. Adenosine: a potential mediator of seizure arrest and postictal refractoriness. \emph{Ann Neurol} 32:618, 1992.
\bibitem{tsodyks1997} Tsodyks M, Markram H. The neural code between neocortical pyramidal neurons depends on neurotransmitter release probability. \emph{PNAS} 94:719, 1997.
\bibitem{wei2014} Wei Y, Ullah G, Schiff SJ. Unification of neuronal spikes, seizures, and spreading depression. \emph{J Neurosci} 34:11733, 2014.
\bibitem{friston2000} Friston KJ, Mechelli A, Turner R, Price CJ. Nonlinear responses in fMRI: the Balloon model, Volterra kernels, and other hemodynamics. \emph{NeuroImage} 12:466, 2000.
\bibitem{ullah2010} Ullah G, Schiff SJ. Assimilating seizure dynamics. \emph{PLoS Comput Biol}, 2010.
\bibitem{wan2000} Wan EA, van der Merwe R. The unscented Kalman filter for nonlinear estimation. IEEE AS-SPCC, 2000.
\bibitem{donoghue2020} Donoghue T et al. Parameterizing neural power spectra into periodic and aperiodic components. \emph{Nat Neurosci} 23:1655, 2020.
\bibitem{gao2017} Gao R, Peterson EJ, Voytek B. Inferring synaptic excitation/inhibition balance from field potentials. \emph{NeuroImage} 158:70, 2017.
\bibitem{scholkmann2013} Scholkmann F, Wolf M. General equation for the differential pathlength factor of the frontal human head depending on wavelength and age. \emph{J Biomed Opt} 18:105004, 2013.
\bibitem{shoeb2009} Shoeb AH. Application of machine learning to epileptic seizure onset detection and treatment. PhD thesis, MIT, 2009 (CHB-MIT database, PhysioNet; Goldberger et al.\ 2000).
\bibitem{bhagubai2025} Bhagubai M, Chatzichristos C, Swinnen L, et al. SeizeIT2: wearable dataset of patients with focal epilepsy. \emph{Scientific Data}, 2025. OpenNeuro ds005873 v1.1.0, doi:10.18112/openneuro.ds005873.v1.1.0 (CC0).
\bibitem{reintjes2025} Reintjes C, Hagenbeck JF, Ballo M, et al. ECG-based detection of epileptic seizures in real-world wearable settings: insights from the SeizeIT2 dataset. \emph{Sensors} 25(24):7687, 2025.
\bibitem{vandecasteele2020} Vandecasteele K et al. Visual seizure annotation and automated seizure detection using behind-the-ear electroencephalographic channels. \emph{Epilepsia} 61:766--775, 2020.
\bibitem{bhagubai2023} Bhagubai M et al. The power of ECG in semi-automated seizure detection in addition to two-channel behind-the-ear EEG. \emph{Bioengineering} 10:491, 2023.
\bibitem{swinnen2021} Swinnen L et al. Accurate detection of typical absence seizures in adults and children using a two-channel EEG wearable behind the ears. \emph{Epilepsia} 62:2741, 2021.
\bibitem{zibrandtsen2017} Zibrandtsen I et al. Ear-EEG detects ictal and interictal abnormalities in focal and generalized epilepsy. \emph{Clin Neurophysiol}, 2017.
\bibitem{jeppesen2019} Jeppesen J et al. Seizure detection based on heart rate variability using a wearable electrocardiography device. \emph{Epilepsia}, 2019.
\bibitem{jeppesen2024} Jeppesen J et al. Detection of seizures with ictal tachycardia using HRV and patient-adaptive logistic regression. \emph{Epileptic Disord}, 2024.
\bibitem{beniczky2018} Beniczky S et al. Automated real-time detection of tonic-clonic seizures using a wearable EMG device. \emph{Neurology} 90, 2018.
\bibitem{onorati2017} Onorati F et al. Multicenter clinical assessment of improved wearable multimodal convulsive seizure detectors (Embrace). \emph{Epilepsia}, 2017; and Embrace real-life case studies (MIT Media Lab).
\bibitem{arends2018} Arends J et al. Multimodal nocturnal seizure detection in a residential care setting (NightWatch). \emph{Neurology}, 2018.
\bibitem{lacuey2018} Lacuey N et al. The incidence and significance of periictal apnea in epileptic seizures. \emph{Epilepsia}, 2018.
\bibitem{hampel2016} Hampel KG et al. Seizure-related modulation of systemic arterial blood pressure in focal epilepsy. \emph{Epilepsia}, 2016.
\bibitem{poh2012} Poh MZ et al. Autonomic changes with seizures correlate with postictal EEG suppression. \emph{Neurology}, 2012.
\bibitem{bubb2023} Bubb C et al. Wearable in-ear pulse oximetry validly measures oxygen saturation between 70\% and 100\%. \emph{Digit Health}, 2023.
\bibitem{inear2022} Accurate detection of heart rate using in-ear photoplethysmography in a clinical setting. \emph{Front Digit Health}, 2022.
\bibitem{prv2021} Accuracy of heart rate variability estimated with reflective wrist-PPG in elderly vascular patients. \emph{Sci Rep}, 2021.
\bibitem{twostage2026} Two-stage seizure detection with wearable single-channel EEG and a gradient-boosting second stage. \emph{IEEE Trans Biomed Eng}, 2026, doi:10.1109/TBME.2026.3660806.
\bibitem{artifact2024} Minimizing artifact-induced false alarms for seizure detection in wearable EEG devices with gradient-boosted tree classifiers. \emph{Sci Rep}, 2024.
\bibitem{phase3ecg2025} Seizure detection using wearable electrocardiogram connected to a smartphone: a phase 3 clinical validation study. \emph{eBioMedicine}, 2025.
\bibitem{pan1985} Pan J, Tompkins WJ. A real-time QRS detection algorithm. \emph{IEEE Trans Biomed Eng} 32:230, 1985.
\bibitem{ilae2014} Fisher RS et al. ILAE official report: a practical clinical definition of epilepsy. \emph{Epilepsia} 55:475, 2014.
\bibitem{ilae2025} ILAE 2025 updated classification of epileptic seizures. \emph{Epilepsia}, 2025.
\bibitem{kwan2000} Kwan P, Brodie MJ. Early identification of refractory epilepsy. \emph{N Engl J Med}, 2000.
\bibitem{who} World Health Organization. Epilepsy fact sheet.
\bibitem{sudep} AAN practice guideline: sudden unexpected death in epilepsy incidence rates and risk factors. \emph{Neurology}, 2017.
\bibitem{temple} Temple (temple.com), Science \& Research; Temple preprints on bioRxiv/medRxiv (time-domain NIRS, transcranial Doppler, age signal), 2026; press coverage as cited in \code{Temple_Technical_Dossier.md}.
\bibitem{goldberger2000} Goldberger AL, Amaral LAN, Glass L, et al. PhysioBank, PhysioToolkit, and PhysioNet: components of a new research resource for complex physiologic signals. \emph{Circulation} 101(23):e215--e220, 2000.
\bibitem{markiewicz2021} Markiewicz CJ, Gorgolewski KJ, Feingold F, et al. The OpenNeuro resource for sharing of neuroscience data. \emph{eLife} 10:e71774, 2021.
\bibitem{pedregosa2011} Pedregosa F, Varoquaux G, Gramfort A, et al. Scikit-learn: machine learning in Python. \emph{J Mach Learn Res} 12:2825--2830, 2011.
\bibitem{ke2017} Ke G, Meng Q, Finley T, et al. LightGBM: a highly efficient gradient boosting decision tree. \emph{Advances in Neural Information Processing Systems} 30, 2017.
\bibitem{paszke2019} Paszke A, Gross S, Massa F, et al. PyTorch: an imperative style, high-performance deep learning library. \emph{Advances in Neural Information Processing Systems} 32, 2019.
\bibitem{lam2015} Lam SK, Pitrou A, Seibert S. Numba: a LLVM-based Python JIT compiler. \emph{Proc.\ LLVM-HPC}, 2015.
\bibitem{kalman1960} Kalman RE. A new approach to linear filtering and prediction problems. \emph{J Basic Eng} 82:35--45, 1960.
\bibitem{beaumont2002} Beaumont MA, Zhang W, Balding DJ. Approximate Bayesian computation in population genetics. \emph{Genetics} 162:2025--2035, 2002.
\bibitem{buxton1998} Buxton RB, Wong EC, Frank LR. Dynamics of blood flow and oxygenation changes during brain activation: the balloon model. \emph{Magn Reson Med} 39:855--864, 1998.
\bibitem{delpy1988} Delpy DT, Cope M, van der Zee P, et al. Estimation of optical pathlength through tissue from direct time of flight measurement. \emph{Phys Med Biol} 33:1433--1442, 1988.
\bibitem{hjorth1970} Hjorth B. EEG analysis based on time domain properties. \emph{Electroencephalogr Clin Neurophysiol} 29:306--310, 1970.
\bibitem{esteller2001} Esteller R, Echauz J, Tcheng T, Litt B, Pless B. Line length: an efficient feature for seizure onset detection. \emph{Proc.\ IEEE EMBS} 2:1707--1710, 2001.
\bibitem{mcsharry2003} McSharry PE, Clifford GD, Tarassenko L, Smith LA. A dynamical model for generating synthetic electrocardiogram signals. \emph{IEEE Trans Biomed Eng} 50:289--294, 2003.
\bibitem{bazett1920} Bazett HC. An analysis of the time-relations of electrocardiograms. \emph{Heart} 7:353--370, 1920.
\bibitem{beniczky2021} Beniczky S, Wiebe S, Jeppesen J, et al. Automated seizure detection using wearable devices: a clinical practice guideline of the International League Against Epilepsy and the International Federation of Clinical Neurophysiology. \emph{Epilepsia} 62:632--646, 2021.
\bibitem{ryvlin2013} Ryvlin P, Nashef L, Lhatoo SD, et al. Incidence and mechanisms of cardiorespiratory arrests in epilepsy monitoring units (MORTEMUS): a retrospective study. \emph{Lancet Neurol} 12:966--977, 2013.
\bibitem{hoppe2007} Hoppe C, Poepel A, Elger CE. Epilepsy: accuracy of patient seizure counts. \emph{Arch Neurol} 64:1595--1599, 2007.
\bibitem{elger2018} Elger CE, Hoppe C. Diagnostic challenges in epilepsy: seizure under-reporting and seizure detection. \emph{Lancet Neurol} 17:279--288, 2018.
\bibitem{detti2020} Detti P, Vatti G, Zabalo Manrique de Lara G. EEG synchronization analysis for seizure prediction: a study on data of noninvasive recordings. \emph{Processes} 8(7):846, 2020. Siena Scalp EEG Database, PhysioNet, doi:10.13026/5d4a-j060.
\bibitem{gu2018} Gu Y, Cleeren E, Dan J, et al. Comparison between scalp EEG and behind-the-ear EEG for development of a wearable seizure detection system for patients with focal epilepsy. \emph{Sensors} 18(1):29, 2018.

\end{thebibliography}
